\VerdierPhysicalPageAnchor{0067}%
\VerdierSemanticAnchor{ch1-2-1-5}
\noindent\textbf{2.1.5.}\quad
Le morphisme $f^\sigma\colon Y^\sigma\longrightarrow Y'^\sigma$ est appelé
le \emph{composant de degré $\sigma$} du morphisme $(f^\sigma)$. Soit
$f\colon Y\longrightarrow Y'$ un morphisme de complexes $n$-uples. Le
composant de degré $\sigma$ de $f$ est noté $f^\sigma$.
\VerdierSemanticAnchor{ch1-2-1-5-composition}
Les morphismes de complexes $n$-uples se composent de manière évidente :
composant par composant.
\VerdierSemanticAnchor{ch1-2-1-5-category}
Les complexes $n$-uples d'objets de $C$ forment une catégorie qui est notée
$\operatorname{comp}^{n}(C)$. Lorsque $n=1$, on utilise la notation simplifiée
$\operatorname{comp}^{1}(C)=\operatorname{comp}(C)$. Les objets de
$\operatorname{comp}(C)$ sont appelés les \emph{complexes simples de $C$} ou
même, lorsqu'aucune confusion n'en résulte, les \emph{complexes de $C$}. Les
catégories $\operatorname{comp}^{n}(C)$ sont additives.

\VerdierSemanticAnchor{ch1-2-1-6}
\noindent\textbf{2.1.6.}\quad
Soit $Y$ un complexe $n$-uple de $C$. On utilise parfois, pour désigner le
composant de degré $\sigma$ de $Y$, la notation $Y_{-\sigma}$, et, pour
désigner la différentielle de degré $\sigma$ de $Y$ dans la direction $i$,
la notation $d^Y_{-\sigma,i}$.
\VerdierSemanticAnchor{ch1-2-1-6-positive-finite}
Cette notation est surtout utilisée lorsque $Y$ est un complexe simple dont
tous les composants de degré positif sont nuls sauf un nombre fini d'entre
eux. Soit de même $f\colon Y\longrightarrow Y'$ un morphisme de complexes
$n$-uples. On utilise parfois la notation $f_{-\sigma}$ pour désigner le
composant de degré $\sigma$ du morphisme $f$.

\VerdierSemanticAnchor{ch1-2-1-7}
\noindent\textbf{2.1.7.}\quad
Soit $\varphi\colon[n]\longrightarrow[m]$ une application.
\VerdierSemanticAnchor{ch1-2-1-7-summable}
Un complexe $n$-uple $Y$ de $C$ est dit \emph{$\varphi$-sommable}, si pour
tout $\tau\in\mathbb Z[m]$, l'ensemble des $\sigma\in\mathbb Z[n]$ tels que
$\varphi(\sigma)=\tau$ et $Y^\sigma$ soit différent de $0$, est fini. Lorsque
$\varphi$ est l'unique application $[n]\longrightarrow[1]$, on dit simplement
\emph{sommable}.
\VerdierSemanticAnchor{ch1-2-1-7-composition}
Soient $[n]\xrightarrow{\varphi}[m]\xrightarrow{\psi}[p]$ deux applications.
Tout complexe $n$-uple $Y$ qui est $\psi\varphi$-sommable est aussi
$\varphi$-sommable. En particulier, tout complexe $n$-uple sommable est
$\varphi$-sommable pour toute application
$[n]\xrightarrow{\varphi}[m]$.
\VerdierSemanticAnchor{ch1-2-1-7-injective}
Lorsque $\varphi$ est une injection, tous les complexes sont
$\varphi$-sommables.
\VerdierSemanticAnchor{ch1-2-1-7-category}
On appelle \emph{catégorie des complexes $n$-uples $\varphi$-sommables}, la
sous-catégorie pleine de $\operatorname{comp}^{n}(C)$ définie par les
complexes $\varphi$-sommables.

\VerdierSemanticAnchor{ch1-2-2}
\subsection{Complexe simple associé}

\VerdierSemanticAnchor{ch1-2-2-1}
\noindent\textbf{2.2.1.}\quad
Pour tout $\sigma\in\mathbb Z[n]$,
$\sigma=\sum_{j\in[n]}l_je_j$, $l_j\in\mathbb Z$, et pour tout $i\in[n]$
posons :

\VerdierSemanticAnchor{ch1-2-2-1-eq1}
\begin{equation}
  \varepsilon(\sigma,i)=(-1)^{\left(\sum_{j<i}l_j\right)}.
  \tag{2.2.1.1}
\end{equation}

\VerdierSemanticAnchor{ch1-2-2-1-sign-map}
La formule (2.2.1.1) définit une application :
\begin{equation*}
  (\sigma,i)\longmapsto\varepsilon(\sigma,i)
\end{equation*}
de $\mathbb Z[n]\times[n]$ dans le groupe des signes $\{-1,1\}$. On vérifie
immédiatement que pour tout $i\in[n]$, l'application
$\sigma\longmapsto\varepsilon(\sigma,i)$ est un homomorphisme de
